\documentclass[11pt,a4paper]{article}

%  XeLaTeX
\usepackage{fontspec}

\usepackage{algorithm}
\usepackage{algorithmic}
\usepackage{amsmath,amssymb,amsthm}
\usepackage{bm}
\usepackage[useregional]{datetime2}
\usepackage{enumitem}
\usepackage{booktabs}
\usepackage{float}
\usepackage{graphicx}
\usepackage{tabularray}
\usepackage{framed}
\usepackage[a4paper, margin=2.5cm]{geometry}
\usepackage{eso-pic}
\usepackage{mathtools}
\usepackage{tikz}
\usetikzlibrary{calc}
\usepackage{xcolor}

\usepackage[backend=biber, style=apa, sorting=nyt]{biblatex}
\addbibresource{D:/github/ENEL445/report/references.bib}
\graphicspath{{D:/github/ENEL445/figs/}}

\usepackage[colorlinks=true, linkcolor=blue, citecolor=blue, urlcolor=blue]{hyperref}
\hypersetup{
  pdfauthor   = {David Ewing},
  pdfsubject  = {\jobname.tex},
  pdfkeywords = {source: \jobname.tex}
}

% Running margin label
\newcommand{\mymarginlabel}{\textcopyright\ 2026 David Ewing dew59@uclive.ac.nz | \DTMnow\ | \jobname.tex}
\AddToShipoutPictureBG{%
  \begin{tikzpicture}[remember picture,overlay]
    \node[rotate=90, anchor=north]
      at ($(current page.north west)+(1.4cm,-13cm)$)
      {\scriptsize\ttfamily \mymarginlabel};
  \end{tikzpicture}%
}

% Custom commands
\newcommand{\E}{\mathbb{E}}
\newcommand{\R}{\mathbb{R}}
\newcommand{\N}{\mathcal{N}}
\newcommand{\ELBO}{\text{ELBO}}
\newcommand{\tr}{\text{tr}}
\newcommand{\sigmoid}{\sigma}
\setlength{\parindent}{0pt}

\title{
Gradient-Based Optimisation for Variational Inference:\\
A Comparative Study Across Linear, Quadratic, and Logistic Regression
}
\author{
  David Ewing (82171165)\\
  ENEL445 --- Applied Engineering Optimisation
}
\date{9 May 2026}

\begin{document}

\maketitle

\begin{abstract}
\setlength{\parindent}{0pt}
\setlength{\parskip}{0.7em}

This report applies the gradient-based optimisation methods from ENEL445 to variational
inference across three Bayesian regression settings: linear, quadratic, and logistic
regression.
The unifying objective is maximising the Evidence Lower Bound (ELBO), which converts
intractable posterior inference into a tractable optimisation problem.

Four optimisation methods are compared in each setting: Coordinate Ascent Variational
Inference (CAVI), gradient ascent with Armijo backtracking, Newton's method with
regularised finite-difference Hessian, and BFGS with Wolfe conditions.
The variational family is $q(\bm{\beta}) = \N(\bm{\mu}, \bm{\Sigma})$ with a Cholesky
reparameterisation to ensure positive definiteness throughout optimisation.

Reference posteriors are obtained from Gibbs samplers.
For the linear and quadratic cases, the conjugate structure admits closed-form Gibbs
updates; the variational family is a Normal--Gamma approximation.
For the logistic case, the Jaakkola--Jordan \parencite{Jaakkola2000} variational bound
restores tractability, and a P\'{o}lya--Gamma Gibbs sampler \parencite{Polson2013}
provides the reference.

The three test cases form a designed progression in inferential difficulty:
the linear case has a conjugate Gaussian likelihood ($n=50$, $p=2$),
the quadratic case adds nonlinearity via feature expansion ($n=100$, $p=3$),
and the logistic case introduces a non-conjugate Bernoulli likelihood ($n=200$, $p=2$).
Posterior variance collapse is mild in the linear case, moderate in the quadratic case,
and minimal in the logistic case.
CAVI is the most efficient method across all three settings.
BFGS is the recommended gradient-based alternative when closed-form CAVI updates are
not available.
\end{abstract}

\newpage
\tableofcontents
\newpage

%─────────────────────────────────────────────────────────────────────────────
\section*{Introduction}
\addcontentsline{toc}{section}{Introduction}
%─────────────────────────────────────────────────────────────────────────────

Bayesian inference provides principled uncertainty quantification for regression models,
but requires computing the posterior distribution
\begin{equation*}
  p(\bm{\beta} \mid \bm{y}) \propto p(\bm{y} \mid \bm{\beta})\,p(\bm{\beta}),
\end{equation*}
which is analytically tractable only for conjugate likelihood--prior pairs.
For most real-world models the posterior integral
\begin{equation*}
  p(\bm{\beta} \mid \bm{y})
  = \frac{p(\bm{y} \mid \bm{\beta})\,p(\bm{\beta})}
         {\displaystyle\int p(\bm{y} \mid \bm{\beta})\,p(\bm{\beta})\, \mathrm{d}\bm{\beta}}
\end{equation*}
cannot be evaluated analytically.
Variational Inference (VI) converts posterior computation into optimisation by finding
the best approximation $q(\bm{\beta})$ from a tractable family, maximising the
Evidence Lower Bound (ELBO):
\begin{equation*}
  \mathcal{L}(\bm{\phi})
  = \E_{q}[\log p(\bm{y}\mid\bm{\beta})]
  - \mathrm{KL}(q \| p)
  \leq \log p(\bm{y}).
\end{equation*}

This project applies the optimisation methods taught in ENEL445 to three Bayesian
regression settings, forming a deliberate progression in inferential difficulty.

\begin{description}
  \item[Case Study 1 --- Linear Regression:]
    Gaussian likelihood with conjugate prior; ELBO has a closed-form gradient.
    Posterior variance collapse is the primary concern.
    Reference posterior from a conjugate Gibbs sampler.
    ($n=50$, $\bm{\beta}_{\mathrm{true}} = (1.0, 2.0)$.)
  \item[Case Study 2 --- Quadratic Regression:]
    Identical Bayesian framework to the linear case, but with feature expansion
    $\bm{x} \mapsto (1, x, x^2)$ so the model remains linear in the coefficient vector.
    Moderate posterior variance collapse.
    ($n=100$, $\bm{\beta}_{\mathrm{true}} = (1.0, 2.0, -1.0)$.)
  \item[Case Study 3 --- Logistic Regression:]
    Non-conjugate Bernoulli likelihood; the Jaakkola--Jordan bound provides a tractable
    ELBO.
    P\'{o}lya--Gamma Gibbs sampler as the reference.
    Posterior variance collapse is minimal.
    ($n=200$, $\bm{\beta}_{\mathrm{true}} = (-1.0, 2.0)$.)
\end{description}

All three studies share the same four optimisation methods (CAVI, gradient ascent,
Newton, BFGS), the same Cholesky reparameterisation for positive-definiteness, and
the same convergence criteria.
The unified methodology allows cross-case comparisons of convergence speed,
computational cost, posterior accuracy, and variance collapse behaviour.

A known failure mode of mean-field VI is \emph{posterior variance collapse}: the
optimised approximation is systematically overconfident, underestimating posterior
uncertainty.
This failure mode is observed and analysed across all three settings.

%─────────────────────────────────────────────────────────────────────────────
\section{Common Methodology}
\label{sec:methodology}
%─────────────────────────────────────────────────────────────────────────────

\subsection{Variational Inference as Optimisation}

The mean-field variational approximation factorises the joint posterior as
$q(\bm{\beta}, \tau_e) = q(\bm{\beta})\,q(\tau_e)$ for the Gaussian regression cases,
or $q(\bm{\beta}) = \N(\bm{\mu}, \bm{\Sigma})$ for the logistic case.
Maximising the ELBO over the variational parameters $\bm{\phi}$ is equivalent to
minimising the KL divergence $\mathrm{KL}(q\|p)$ between the approximation and the
true posterior.

\subsection{Constraint Handling: Cholesky Reparameterisation}

The covariance matrix $\bm{\Sigma}$ must remain symmetric positive definite at every
iteration.
The set of symmetric positive definite matrices has a curved boundary: a standard
Euclidean gradient step can cross outside it, producing an invalid covariance.
The solution is to parameterise over the Cholesky factor $\bm{L}$, where
$\bm{\Sigma} = \bm{L}\bm{L}^T$ and $\bm{L}$ is lower-triangular with positive diagonal
entries.
The diagonal entries are reparameterised as $L_{ii} = e^{\tilde{\ell}_{ii}}$ to allow
unconstrained optimisation.
This converts the constraint into a smooth bijection and ensures every gradient step
produces a valid covariance.

\subsection{Optimisation Methods}

All four methods share the same objective (the ELBO), the same reparameterisation,
and the same convergence criteria.

\begin{description}
  \item[CAVI:] Coordinate Ascent Variational Inference.
    Closed-form coordinate updates derived from the exponential family structure.
    Each update exactly maximises the ELBO over one block of parameters while holding
    the others fixed.
    Most efficient method; requires conjugate structure.
  \item[Gradient ascent:] First-order method.
    Search direction is the gradient $\bm{d}^{(t)} = \nabla\mathcal{L}(\bm{\phi}^{(t)})$.
    Step size determined by Armijo backtracking (sufficient increase condition).
  \item[Newton's method:] Second-order method.
    Search direction is $\bm{d}^{(t)} = -\bm{H}^{-1}\nabla\mathcal{L}$,
    where $\bm{H}$ is the Hessian approximated by finite differences (25 evaluations
    per iteration).
    Fast convergence near the optimum; expensive per iteration.
  \item[BFGS:] Quasi-Newton method.
    Builds a curvature approximation from gradient differences, requiring only
    gradient evaluations.
    Step size determined by Wolfe conditions to maintain positive definiteness of
    the approximate inverse Hessian ($\bm{y}_t^T\bm{s}_t > 0$).
    Best practical compromise between convergence speed and cost.
\end{description}

For all line-search methods, the search direction must be an ascent direction:
$\nabla\mathcal{L}(\bm{\phi}^{(t)})^T\bm{d}^{(t)} > 0$.

\subsection{Convergence Criteria}

Iteration terminates when any of the following conditions is satisfied:
\begin{enumerate}[label=(\roman*)]
  \item Gradient norm: $\lVert\nabla\mathcal{L}\rVert < \varepsilon_g$.
  \item Relative ELBO change: $|\Delta\mathcal{L}|/|\mathcal{L}| < \varepsilon_r$.
  \item Parameter change: $\lVert\Delta\bm{\phi}\rVert < \varepsilon_\phi$.
  \item Maximum iterations: $t \geq t_{\max}$.
\end{enumerate}
CAVI additionally monitors the maximum relative change in variational parameters.

\subsection{Reference Samplers}

Two types of Gibbs sampler are used as gold-standard references.

\paragraph{Conjugate Gibbs (linear and quadratic cases).}
The Normal--Gamma conjugate structure of the Bayesian regression model (Gaussian
likelihood with Gaussian prior on $\bm{\beta}$ and Gamma prior on precision $\tau_e$)
admits closed-form Gibbs updates:
\begin{align}
  \bm{\beta} \mid \tau_e, \bm{y} &\sim \N(\bm{\mu}_n, \bm{\Sigma}_n), \\
  \tau_e \mid \bm{\beta}, \bm{y} &\sim \mathrm{Gamma}(a_n, b_n),
\end{align}
where $\bm{\Sigma}_n = (\tau_e\bm{X}^T\bm{X} + \lambda_\beta\bm{I})^{-1}$ and
$\bm{\mu}_n = \tau_e\bm{\Sigma}_n\bm{X}^T\bm{y}$.

\paragraph{P\'{o}lya--Gamma Gibbs (logistic case).}
The P\'{o}lya--Gamma (PG) sampler of \textcite{Polson2013} augments the logistic model
with latent variables $\omega_i \sim \mathrm{PG}(1, x_i^T\bm{\beta})$, rendering
$\bm{\beta}$ conditionally Gaussian:
\begin{align}
  \bm{\beta} \mid \bm{\omega}, \bm{y}
    &\sim \N\!\left(\bm{\Sigma}_\omega(\bm{X}^T\bm{\kappa}),\; \bm{\Sigma}_\omega\right),
  \quad \bm{\Sigma}_\omega = (\bm{X}^T\bm{\Omega}\bm{X} + \bm{\Sigma}_0^{-1})^{-1},
\end{align}
where $\bm{\Omega} = \mathrm{diag}(\omega_i)$ and $\kappa_i = y_i - 1/2$.
The PG draw requires $T$ Gamma variates per latent variable; the truncation $T=50$
is used throughout.

\subsection{Common Notation}

\begin{center}
\begin{tabular}{@{}ll@{}}
\toprule
Symbol & Meaning \\
\midrule
\multicolumn{2}{@{}l}{\textbf{Bayesian model}} \\
$\bm{X}$ & Design matrix ($n \times p$), rows $x_i^T$ \\
$\bm{y}$ & Response vector \\
$\bm{\beta}$ & Regression coefficient vector \\
$\tau_e$ & Observation-noise precision (Gaussian cases) \\
$n,\,p$ & Number of observations; number of predictors \\
$\lambda_\beta$ & Prior precision on $\bm{\beta}$ \\
\midrule
\multicolumn{2}{@{}l}{\textbf{Variational approximation}} \\
$q(\bm{\beta})$ & Variational Gaussian approximation \\
$\bm{\mu},\,\bm{\Sigma}$ & Mean and covariance of $q(\bm{\beta})$ \\
$\bm{L}$ & Cholesky factor, $\bm{\Sigma}=\bm{L}\bm{L}^T$ \\
$\mathcal{L}(\bm{\phi})$ & Evidence Lower Bound (ELBO) \\
\midrule
\multicolumn{2}{@{}l}{\textbf{Optimisation}} \\
$\bm{d}^{(t)}$ & Search direction at iteration $t$ \\
$\alpha_t$ & Step size at iteration $t$ \\
$\bm{s}_t,\,\bm{y}_t$ & Parameter-change and gradient-difference vectors (BFGS) \\
\midrule
\multicolumn{2}{@{}l}{\textbf{Logistic-specific}} \\
$\sigma(\cdot)$ & Logistic sigmoid, $\sigma(\psi)=(1+e^{-\psi})^{-1}$ \\
$\xi_i$ & Jaakkola--Jordan tightness parameter \\
$\omega_i$ & P\'{o}lya--Gamma latent variable \\
$\kappa_i$ & Pseudo-response $\kappa_i = y_i - 1/2$ \\
$\mathrm{ESS}$ & Effective sample size \\
\midrule
\multicolumn{2}{@{}l}{\textbf{Lecture correspondence (Prof.\ R.\,S.\ Smith, ENEL445 2026)}} \\
$J(\theta)$ & Cost function $= \mathcal{L}(\bm{\phi})$ (ELBO) \\
$\Phi$ & Regressor matrix $= \bm{X}$ \\
\bottomrule
\end{tabular}
\end{center}

%─────────────────────────────────────────────────────────────────────────────
\section{Case Study 1: Bayesian Linear Regression}
\label{sec:linear}
%─────────────────────────────────────────────────────────────────────────────

\subsection{Model}

The Bayesian linear regression model is
\begin{align}
  \bm{y} &\sim \N(\bm{X}\bm{\beta},\; \tau_e^{-1}\bm{I}_n), \\
  \bm{\beta} &\sim \N(\bm{0},\; \lambda_\beta^{-1}\bm{I}_p), \quad \lambda_\beta = 0.1, \\
  \tau_e &\sim \mathrm{Gamma}(\alpha_e, \gamma_e).
\end{align}
The mean-field variational approximation is
$q(\bm{\beta},\tau_e) = \N(\bm{\mu}_\beta,\bm{\Sigma}_\beta)\cdot\mathrm{Gamma}(a_e,b_e)$.

The test case uses $n=50$ observations with $x_i \sim \mathrm{Uniform}(-2,2)$ and
true parameters $\bm{\beta}_{\mathrm{true}} = (1.0,\; 2.0)^T$ ($p=2$).
Data are generated with noise standard deviation $\sigma=0.5$.

\subsection{ELBO}

Under the Normal--Gamma variational family, the ELBO is
\begin{align}
  \mathcal{L}(\bm{\phi})
  &= \E_q[\log p(\bm{y}\mid\bm{\beta},\tau_e)]
   + \E_q[\log p(\bm{\beta})]
   + \E_q[\log p(\tau_e)]
   - \E_q[\log q(\bm{\beta})]
   - \E_q[\log q(\tau_e)].
\end{align}
Each term has a closed-form expression:
\begin{align}
  \E_q[\log p(\bm{y}\mid\bm{\beta},\tau_e)]
  &= \frac{n}{2}\bigl(\psi(a_e)-\log b_e\bigr)
    - \frac{a_e}{2b_e}\bigl(
        \lVert\bm{y}-\bm{X}\bm{\mu}_\beta\rVert^2
        + \tr(\bm{X}\bm{\Sigma}_\beta\bm{X}^T)
      \bigr)
    - \frac{n}{2}\log(2\pi),\\
  \mathrm{KL}(q(\bm{\beta})\|p(\bm{\beta}))
  &= \tfrac{1}{2}\bigl[
      \lambda_\beta\,\tr(\bm{\Sigma}_\beta)
      + \lambda_\beta\lVert\bm{\mu}_\beta\rVert^2
      - p
      - \log\det(\lambda_\beta\bm{\Sigma}_\beta)
    \bigr],
\end{align}
with analogous terms for the Gamma factor.
Full derivations are in the individual linear report (Appendix~\ref{app:linear-ref}).

\subsection{Results}

[dedbug - ./linear\_run.ipynb}.\\
* linear\_*.png}.\\
* linear\_diagnostics.csv}.]

%─────────────────────────────────────────────────────────────────────────────
\section{Case Study 2: Bayesian Quadratic Regression}
\label{sec:quadratic}
%─────────────────────────────────────────────────────────────────────────────

\subsection{Model}

The quadratic regression case uses the same Bayesian framework as the linear case,
but with a feature-expanded design matrix.
The raw covariate $x_i \sim \mathrm{Uniform}(-2,2)$ is mapped to
$\tilde{x}_i = (1,\; x_i,\; x_i^2)^T$, so the model is
\begin{equation}
  y_i = \beta_0 + \beta_1 x_i + \beta_2 x_i^2 + \varepsilon_i,
  \quad \varepsilon_i \sim \N(0, \tau_e^{-1}).
\end{equation}
The model is linear in $\bm{\beta} = (\beta_0, \beta_1, \beta_2)^T$, so the same
Normal--Gamma ELBO and CAVI updates apply with $p=3$.

The test case uses $n=100$ observations and true parameters
$\bm{\beta}_{\mathrm{true}} = (1.0,\; 2.0,\; -1.0)^T$.
The quadratic term $\beta_2 = -1.0$ introduces a downward curvature that tests whether
the optimisers correctly attribute variance to the additional parameter.

\subsection{ELBO}

The ELBO is identical in form to the linear case (Section~\ref{sec:linear}) with
$p=3$ and the expanded design matrix $\tilde{\bm{X}}$.
There is no new approximation: the feature expansion is a deterministic preprocessing
step applied before inference.

\subsection{Results}

[deebug - ./quadratic\_run.ipynb}.\\
* quadratic\_*.png}.\\
* quadratic\_diagnostics.csv}.]

%─────────────────────────────────────────────────────────────────────────────
\section{Case Study 3: Bayesian Logistic Regression}
\label{sec:logistic}
%─────────────────────────────────────────────────────────────────────────────

\subsection{Model}

The Bayesian logistic regression model is
\begin{align}
  y_i \mid \bm{\beta} &\sim \mathrm{Bernoulli}(\sigma(x_i^T\bm{\beta})),
  \quad i = 1,\ldots,n, \\
  \bm{\beta} &\sim \N(\bm{0},\; \lambda^{-1}\bm{I}_p),
  \quad \lambda^{-1} = 10,
\end{align}
where $\sigma(\psi) = (1+e^{-\psi})^{-1}$ is the logistic sigmoid.
The logistic likelihood is not conjugate to the Gaussian prior, so the posterior is not
available in closed form.

The test case uses $n=200$ observations with $x_i \sim \mathrm{Uniform}(-2,2)$ and
true parameters $\bm{\beta}_{\mathrm{true}} = (-1.0,\; 2.0)^T$ ($p=2$).

\subsection{Jaakkola--Jordan ELBO}

The Jaakkola--Jordan bound \parencite{Jaakkola2000} introduces local variational
parameters $\xi_i > 0$ to replace the intractable log-likelihood term
$\E_q[\log\sigma(x_i^T\bm{\beta})]$ with a tight quadratic lower bound.
At the optimal tightness $\xi_i^* = \sqrt{\E_q[(x_i^T\bm{\beta})^2]}$, the ELBO is
\begin{equation}
  \mathcal{L}(\bm{\theta})
  = \sum_{i=1}^{n}
    \Bigl[(y_i - \tfrac{1}{2})\,x_i^T\bm{\mu}
          - \log 2 - \log\cosh\!\tfrac{\xi_i^*}{2}\Bigr]
  - \mathrm{KL}(q\|p).
\end{equation}
The gradient and CAVI updates for this ELBO are derived in \textcite{Jaakkola2000};
the implementation uses Cholesky reparameterisation with 5 unconstrained parameters.

\subsection{Results}

We generated synthetic data using a known logistic regression model with parameters
$\beta_0 = -1.0$ and $\beta_1 = 2.0$.
These are the data-generating parameters --- they define the process that produced the
observations and represent what we are trying to recover.

Inference was conducted as if the true values were unknown.
The Gibbs sampler uses only two inputs: the 200 binary observations and the weakly
informative prior $\bm{\beta} \sim \N(\bm{0}, 10\bm{I})$.
It never sees the true values.

The posterior distribution is the result of combining what the data say (the likelihood)
with what we believed before seeing the data (the prior).
The posterior mean sits at approximately $(-1.6,\; 2.6)$, offset from the true values,
because with only 200 observations the likelihood is not strong enough to overcome the
prior completely.

The true values appear on the plots only as a benchmark --- to verify that the posterior
is plausible.
The fact that the true values fall within the posterior distribution confirms the sampler
is behaving correctly.

\begin{figure}[H]
  \centering
  \includegraphics[width=0.85\textwidth]{logistic_elbo_convergence.png}
  \caption{ELBO convergence for all four optimisation methods on the logistic test case.
           All methods converge to the same final value within 50 iterations.}
  \label{fig:logistic-elbo}
\end{figure}

\begin{figure}[H]
  \centering
  \includegraphics[width=0.85\textwidth]{logistic_vb_gibbs_beta0.png}
  \caption{Posterior comparison for $\beta_0$: variational approximation versus PG Gibbs
           reference density.
           The posterior mean is offset from the true value $\beta_0=-1.0$ because
           200 observations are insufficient to fully override the prior.}
  \label{fig:logistic-posterior}
\end{figure}

\begin{figure}[H]
  \centering
  \includegraphics[width=0.85\textwidth]{logistic_gibbs_ess.png}
  \caption{Effective sample size (ESS) per parameter for the PG Gibbs sampler.
           ESS of approximately 250--280 out of 2000 samples indicates moderate
           autocorrelation ($\approx 8\times$ cost penalty per effective sample).}
  \label{fig:logistic-ess}
\end{figure}

%─────────────────────────────────────────────────────────────────────────────
\section{Discussion}
\label{sec:discussion}
%─────────────────────────────────────────────────────────────────────────────

\subsection{Cross-Case Comparison}

Table~\ref{tab:comparison} summarises the key structural and empirical differences
across the three test cases.
All timing figures are from the same machine (Alienware 17~R5, Intel i7-8750H) but
are not controlled comparisons: model complexity, parameter count, and Gibbs overhead
differ across cases.
The fair within-notebook comparison is each VI method versus its own Gibbs reference.

\begin{table}[H]
\centering
\caption{Comparison of the three case studies.}
\label{tab:comparison}
\begin{tabular}{@{}llll@{}}
\toprule
Property & Linear & Quadratic & Logistic \\
\midrule
Likelihood & Gaussian & Gaussian & Bernoulli \\
True $\bm{\beta}$ & $(1.0, 2.0)$ & $(1.0, 2.0, -1.0)$ & $(-1.0, 2.0)$ \\
Predictors $p$ & 2 & 3 & 2 \\
Observations $n$ & 50 & 100 & 200 \\
Conjugacy & Yes & Yes & No (JJ bound) \\
ELBO form & Normal--Gamma & Normal--Gamma & Jaakkola--Jordan \\
Reference sampler & Conjugate Gibbs & Conjugate Gibbs & PG Gibbs \\
Gibbs cost & Low & Low & High ($n \times T$ Gammas) \\
Variational parameters & $p + p(p+1)/2 + 2$ & $p + p(p+1)/2 + 2$ & $2p + 1$ (Cholesky) \\
Variance collapse & Mild & Moderate & Minimal \\
CAVI updates & Closed form & Closed form & Closed form (via JJ) \\
\bottomrule
\end{tabular}
\end{table}

\subsection{Posterior Variance Collapse}

Posterior variance collapse --- the tendency of mean-field VI to underestimate posterior
uncertainty --- varies systematically across the three cases.

In the linear and quadratic Gaussian cases, the mean-field factorisation
$q(\bm{\beta})\,q(\tau_e)$ removes true posterior correlations between $\bm{\beta}$ and
$\tau_e$, squeezing the marginal approximations.
The effect is mild for the linear case (small $p=2$ and small $n=50$) and more pronounced
for the quadratic case (larger $p=3$ and $n=100$).

In the logistic case with the Jaakkola--Jordan bound, variance collapse is minimal.
This is because the JJ bound is tight at the CAVI fixed point: when
$\xi_i^* = \sqrt{\E_q[(x_i^T\bm{\beta})^2]}$, the bound equals the true expectation,
and the ELBO objective penalises over-concentration.
The result is that CAVI and all gradient-based methods recover posterior standard
deviations accurately relative to the PG Gibbs reference.

\subsection{Optimisation Behaviour}

CAVI is the most efficient method across all three settings, because it exploits
closed-form coordinate updates.
For the linear and quadratic cases, the conjugate Normal--Gamma structure enables
exact updates; for the logistic case, the JJ bound restores the same conjugate structure.

BFGS converges rapidly in all three cases, providing an effective gradient-based
alternative when closed-form CAVI updates are not available (as would be the case for
non-standard likelihoods or non-conjugate priors).

Newton's method is the most expensive per-iteration due to the finite-difference
Hessian, but converges in the fewest iterations near the optimum.
For the problem sizes here ($p \leq 3$), the finite-difference overhead is manageable.

%─────────────────────────────────────────────────────────────────────────────
\section*{Conclusion}
\addcontentsline{toc}{section}{Conclusion}
%─────────────────────────────────────────────────────────────────────────────

This project has applied gradient-based optimisation methods to variational inference
across three Bayesian regression settings of increasing difficulty.
The unifying principle is that ELBO maximisation transforms intractable posterior
inference into a tractable optimisation problem, directly amenable to the methods
taught in ENEL445.

Across all three case studies, CAVI is the most efficient method, converging rapidly
via closed-form coordinate updates.
BFGS is the recommended gradient-based alternative, achieving quasi-Newton convergence
speed with only gradient evaluations and Wolfe-condition step control.

The three cases expose a clear progression:
\begin{enumerate}
  \item The linear case establishes the baseline: conjugate Gibbs reference, mild
    variance collapse, all methods converge cleanly.
  \item The quadratic case raises the parameter count and reveals that moderate variance
    collapse can persist even when the ELBO has converged.
  \item The logistic case introduces a non-conjugate computational bottleneck, resolved
    by the Jaakkola--Jordan bound, and demonstrates that variance collapse can be
    avoided entirely when the bound is tight at the CAVI fixed point.
\end{enumerate}

The P\'{o}lya--Gamma Gibbs sampler provides an accurate reference for the logistic case
but is computationally expensive ($n \times T$ Gamma draws per iteration).
For production settings with $n \gg 200$, the VI methods offer a substantial
computational advantage.

\printbibliography

%─────────────────────────────────────────────────────────────────────────────
\appendix
%─────────────────────────────────────────────────────────────────────────────

\section{Case Study Summary Table}
\label{app:comparison-table}

The following table expands on Table~\ref{tab:comparison} with additional
optimisation-specific and sampler-specific properties.

\begin{table}[H]
\centering
\caption{Extended case study comparison.}
\begin{tabular}{@{}llll@{}}
\toprule
Property & Linear & Quadratic & Logistic \\
\midrule
\multicolumn{4}{@{}l}{\textbf{Model}} \\
Likelihood & Gaussian & Gaussian & Bernoulli \\
True $\bm{\beta}$ & $(1.0, 2.0)$ & $(1.0, 2.0, -1.0)$ & $(-1.0, 2.0)$ \\
Predictors $p$ & 2 & 3 & 2 \\
Observations $n$ & 50 & 100 & 200 \\
Conjugacy & Yes & Yes & No \\
\midrule
\multicolumn{4}{@{}l}{\textbf{Variational approximation}} \\
Variational family & Normal--Gamma & Normal--Gamma & Gaussian only \\
ELBO form & Closed form & Closed form & Jaakkola--Jordan \\
VI parameters $d$ & 8 & 11 & 5 \\
Variance collapse & Mild & Moderate & Minimal \\
\midrule
\multicolumn{4}{@{}l}{\textbf{Reference sampler}} \\
Sampler type & Conjugate Gibbs & Conjugate Gibbs & PG Gibbs \\
Cost per sweep & $O(p^3)$ & $O(p^3)$ & $O(n \cdot T)$ \\
ESS (typical) & High & High & $\approx 250$--$280$ \\
\midrule
\multicolumn{4}{@{}l}{\textbf{Optimisation}} \\
CAVI convergence & $<50$ iterations & $<50$ iterations & $<50$ iterations \\
BFGS convergence & Fast & Fast & Fast \\
Recommended method & CAVI & CAVI & CAVI \\
\bottomrule
\end{tabular}
\end{table}

The text associated with the logistic case for presentation

\begin{framed}
\noindent
We generated synthetic data using a known logistic regression model with parameters
$\beta_0 = -1.0$ and $\beta_1 = 2.0$.
These are the data-generating parameters: they define the process that produced the
observations and represent what we are trying to recover.

We then treated inference as if the true values were unknown.
The Gibbs sampler uses only two inputs: the 200 binary observations and a weakly
informative prior $\bm{\beta} \sim \N(\bm{0}, 10\bm{I})$.
It never sees the true values.

The posterior distribution is the result of combining what the data say (the likelihood)
with what we believed before seeing the data (the prior).
The posterior mean sits at approximately $(-1.6, 2.6)$, offset from the true values,
because with only 200 observations the likelihood is not strong enough to overcome
the prior completely.

The true values appear on the plots only as a benchmark --- to verify that the posterior
is plausible.
The fact that the true values fall within the posterior distribution confirms the sampler
is behaving correctly.
\end{framed}

%─────────────────────────────────────────────────────────────────────────────
\section{Individual Report References}
\label{app:individual-refs}
%─────────────────────────────────────────────────────────────────────────────

The three individual case study reports contain full derivations, all figures, and
complete code listings.

\begin{description}
  \item[Linear regression:] \texttt{report/20260505D-LINEAR-REPORT.tex}.
    \label{app:linear-ref}
    Full ELBO derivation, gradient verification, ELBO landscape, all convergence
    figures, posterior comparison figures, and variance collapse analysis.
  \item[Quadratic regression:] \texttt{report/20260506B-QUADRATIC-REPORT.tex}.
    \label{app:quad-ref}
    Identical framework to the linear case with $p=3$ and feature expansion.
  \item[Logistic regression:] \texttt{report/20260509A-LOGISTIC-REPORT.tex}.
    \label{app:logistic-ref}
    Full Jaakkola--Jordan derivation, P\'{o}lya--Gamma sampler, ESS diagnostics,
    and all 20 result figures.
\end{description}

%─────────────────────────────────────────────────────────────────────────────
\section{Code Structure}
\label{app:code}
%─────────────────────────────────────────────────────────────────────────────

All code is in the \texttt{python/} directory.
Each notebook follows a five-stage structure.

\begin{description}
  \item[Stage 1:] Data generation. Saves \texttt{results/data/<case>\_data.parquet}.
  \item[Stage 2:] ELBO evaluation and gradient verification (finite-difference check).
  \item[Stage 3:] Optimisation (CAVI, gradient ascent, Newton, BFGS).
    Saves \texttt{results/data/<case>\_opt.parquet}.
  \item[Stage 4:] Reference Gibbs sampler with ESS diagnostics.
    Saves \texttt{results/data/<case>\_gibbs.parquet}.
  \item[Stage 5:] Diagnostics and comparison figures.
    Saves \texttt{results/tables/<case>\_diagnostics.csv}.
\end{description}

\begin{center}
\begin{tabular}{@{}lll@{}}
\toprule
Case & Notebook & Figure prefix \\
\midrule
Linear     & \texttt{python/linear\_run.ipynb}     & \texttt{linear\_} \\
Quadratic  & \texttt{python/quadratic\_run.ipynb}  & \texttt{quadratic\_} \\
Logistic   & \texttt{python/logistic\_run.ipynb}   & \texttt{logistic\_} \\
\bottomrule
\end{tabular}
\end{center}

Core algorithms are in \texttt{python/vb\_algorithms\_py.py};
utilities (archive, plot, I/O) are in \texttt{python/vb\_utils\_py.py}.

All four reports are available online at:
\begin{center}
  \url{https://ENEL445.IndustrialOrigami.ai}
\end{center}

\end{document}
